\section{高考考频分析}

\subsection{近五年考点分布（2021--2025 新课标Ⅰ卷）}

\begin{center}
\begin{tabular}{c|c|c|c|c|c}
\textbf{知识模块} & \textbf{2021} & \textbf{2022} & \textbf{2023} & \textbf{2024} & \textbf{2025} \\ \hline
分子与细胞 & 15 & 18 & 12 & 15 & 18 \\
遗传与进化 & 22 & 20 & 22 & 25 & 20 \\
稳态与调节 & 22 & 25 & 22 & 18 & 22 \\
生物与环境 & 12 & 10 & 12 & 15 & 10 \\
生物技术与工程 & 15 & 15 & 15 & 15 & 15 \\
实验与探究 & 14 & 12 & 17 & 12 & 15 \\ \hline
\textbf{总计} & 100 & 100 & 100 & 100 & 100
\end{tabular}
\end{center}

\subsection{必考/常考/轮考模块}

\[
\begin{array}{c|c|c|c}
\text{类别} & \text{模块} & \text{年均分值} & \text{考查形式} \\ \hline
\textbf{必考} & 遗传基本定律 & 8--12 分 & 非选择题（遗传推理+计算） \\
\textbf{必考} & 神经体液调节 & 8--12 分 & 非选择题（调节过程分析） \\
\textbf{必考} & 细胞代谢（光合+呼吸） & 6--10 分 & 非选择题（过程图解+实验） \\
\textbf{必考} & 生态系统 & 6--8 分 & 非选择题（能量流动+物质循环） \\
\textbf{常考} & 基因表达 & 4--6 分 & 选择题+填空 \\
\textbf{常考} & 免疫调节 & 4--6 分 & 选择题+填空 \\
\textbf{常考} & 种群与群落 & 4--6 分 & 选择题+填空 \\
\textbf{常考} & 细胞分裂 & 4--6 分 & 选择题（图像辨析） \\
\textbf{轮考} & 植物激素 & 2--4 分 & 选择题 \\
\textbf{轮考} & 进化 & 2--4 分 & 选择题 \\
\textbf{轮考} & 内环境稳态 & 2--4 分 & 选择题
\end{array}
\]

\subsection{题型分布特征}

\textbf{选择题分布}（16 题 × 2.5 分 = 40 分）：
\begin{itemize}
  \item 第 1--4 题：分子与细胞基础（化合物、细胞结构、酶、ATP）
  \item 第 5--8 题：遗传与进化基础（DNA、基因表达、变异育种）
  \item 第 9--12 题：稳态与调节（神经、体液、免疫、植物激素）
  \item 第 13--16 题：生态+实验+选修综合（图表信息题为主）
\end{itemize}

\textbf{非选择题分布}（5 题共 60 分）：
\[
\begin{array}{c|c|c}
\text{题号} & \text{分值} & \text{常考内容} \\ \hline
17 & 8--10 分 & 细胞代谢（光合作用/细胞呼吸） \\
18 & 10--12 分 & 调节（神经/体液/免疫） \\
19 & 8--10 分 & 生态（能量流动/物质循环/种群） \\
20 & 12--16 分 & 遗传（定律/系谱图/变异） \\
21 & 10--15 分 & 选修（生物技术/现代生物科技）
\end{array}
\]

\subsection{命题趋势分析}

\begin{enumerate}
  \item \textbf{情境化程度加深}：近年试题大量结合科研文献、生活实际、农业生产等真实情境
  \item \textbf{图表信息题占比增加}：曲线图、柱状图、表格、流程图等形式考查信息提取能力
  \item \textbf{实验探究能力要求提高}：实验设计、结果预测、结论分析类题目增多
  \item \textbf{长句表达考查增加}："原因说明""依据""判断理由"等需要完整作答
  \item \textbf{选修内容交叉融合}：选修与必修知识的综合考查趋势明显
  \item \textbf{遗传计算回归基础}：近年遗传计算难度有所降低，更注重基本原理与推理过程
\end{enumerate}

\subsection{复习优先级建议}

\[
\begin{array}{c|c|c}
\text{优先级} & \text{模块} & \text{理由} \\ \hline
\textbf{第一梯队} & 代谢、遗传、调节 & 分值最高，非选择题必考 \\
\textbf{第二梯队} & 生态、基因表达 & 常考且易拿分 \\
\textbf{第三梯队} & 细胞分裂、变异育种、免疫 & 选择题高频考点 \\
\textbf{第四梯队} & 进化、植物激素、内环境 & 轮考，了解基础即可
\end{array}
\]

\textbf{目标 60 分}：重点突破第一梯队基础知识 + 第二梯队全部 + 选修模板。

\textbf{目标 80 分}：在第一梯队基础上强化实验设计 + 长句表达 + 遗传推理。

\textbf{目标 90+}：全模块无死角，重点突破图表信息题和创新情境题。
